\section{观测模型与定位目标}\label{sec:model}
\subsection{同步单天线站的宽带观测}
在静态直达传播条件下，源位置决定了信号到各站的相对幅度与时延。设$L$个单天线站的已知位置为$\bm r_l\in\R^2$，$l=1,\ldots,L$；$K$个源的未知位置为$\bm p_k\in\Omega\subset\R^2$，$k=1,\ldots,K$，其中$\Omega$为二维搜索区域。各站时间、频率同步，接收链路相对幅相响应已标定；发射天线各向一致，接收天线具有相同响应并满足极化匹配。本文先研究已知源数的定位，实验取$L=4$、$K=2$。距离、时延与相对幅度分别为
\begin{equation}
 d_{lk}=\|\bm p_k-\bm r_l\|_2,\qquad
 \tau_{lk}=d_{lk}/c_0,\qquad \beta_{lk}=d_{lk}^{-1},
 \label{eq:distance}
\end{equation}
其中$c_0$为传播速度，$d_{lk}$和$\tau_{lk}$的单位分别为m和s。

幅度模型来自自由空间路径损耗。对于远场、极化匹配的理想各向同性收发天线，接收功率$P_{\mathrm r}$与发射功率$P_{\mathrm t}$满足\cite{ITURP5252024}
\begin{equation}
 \frac{P_{\mathrm r}}{P_{\mathrm t}}
 =\left(\frac{\lambda}{4\pi d}\right)^2,\qquad
 a_{\mathrm{amp}}(d)=\sqrt{\frac{P_{\mathrm r}}{P_{\mathrm t}}}
 =\frac{\lambda}{4\pi d} .\label{eq:freespace}
\end{equation}
这里$\lambda$为工作波长，$d$为距离，$a_{\mathrm{amp}}$为幅度比。式\eqref{eq:freespace}说明，功率按$1/d^2$衰减，幅度按$1/d$衰减。工作波长和天线响应固定时，将公共因子$\lambda/(4\pi)$并入未知源波形，就得到$\beta_{lk}=1/d_{lk}$。定位器因此能利用站间相对幅度，而无需知道发射功率。本文用该自由空间规律描述幅度，用基带时延描述时间与相位关系，不单独建模射频载波相位。

令$s_k(t)$为已包含公共幅度因子和未知发射功率的第$k$个源基带波形。第$l$站以采样率$f_s$采集$N$点，接收信号为
\begin{equation}
 x_l[n]=\sum_{k=1}^{K}\beta_{lk}s_k(n/f_s-\tau_{lk})+w_l[n],
 \quad n=0,\ldots,N-1 .\label{eq:time}
\end{equation}
其中$x_l[n]$为复数IQ样本，$n$为采样索引。源波形互相独立且未知，在仿真中采用独立、零均值的随机符号产生。噪声$w_l[n]\sim\mathcal{CN}(0,\sigma_t^2)$为圆对称复高斯白噪声，$\sigma_t^2=\E|w_l[n]|^2$，在站间及采样点间独立，并与源信号独立。给定候选位置后，其到各站的相对幅度与时延均由式\eqref{eq:distance}确定。

为利用宽带信号的时延信息，对每站记录作离散傅里叶变换（DFT）：
\begin{equation}
 X_l(f_m)=\frac1N\sum_{n=0}^{N-1}x_l[n]
 \exp(-\mathrm j2\pi mn/N),\qquad f_m=mf_s/N .\label{eq:dft}
\end{equation}
$m$采用移位后的有符号频率索引，$\mathrm j=\sqrt{-1}$。在DFT频率轴上保留公共处理频带内的频点，组成集合$\mathcal F$，其大小记为$F=|\mathcal F|$。以$\bm x_f=[X_1(f),\ldots,X_L(f)]^T\in\C^L$收集同一频点的各站观测，则宽带频域模型写为\cite{ChenHudsonYao2002}
\begin{align}
 \bm x_f&=\bm A_f(\bm P)\bm s_f+\bm n_f,\quad f\in\mathcal F,\label{eq:frequency}\\
 [\bm a_f(\bm p)]_l&=\frac{1}{d_l(\bm p)}
       \exp[-\mathrm j2\pi f d_l(\bm p)/c_0],\label{eq:response}\\
 \bm A_f(\bm P)&=[\bm a_f(\bm p_1),\ldots,\bm a_f(\bm p_K)],\quad
 \bm P=[\bm p_1,\ldots,\bm p_K].\nonumber
\end{align}
其中$\bm s_f=[S_1(f),\ldots,S_K(f)]^T\in\C^K$为未知源频域系数，$\bm n_f\in\C^L$为噪声，$\bm a_f(\bm p)\in\C^L$为候选位置的跨站响应，$\bm A_f(\bm P)\in\C^{L\times K}$为各源响应组成的矩阵，$\bm P\in\R^{2\times K}$收集源坐标。式\eqref{eq:response}的相位项来自时移的傅里叶变换性质，$f$表示基带频率。全频带需要逐频计算响应，不能以单一频率的导向向量代替。

上述频域模型采用共同观测窗的时延近似，忽略各站有限记录在窗边缘的差异。完整DFT的白噪声满足$\E[\bm n_f\bm n_f^H]=(\sigma_t^2/N)\bm I_L$。去除带外频点后，保留频点的噪声方差不变，DFT归一化仍使用原始点数$N$。所有方法采用这一处理，不对各站混合IQ作额外RMS归一化。

计算中也可使用单位范数响应$\bm q_f(\bm p)=\bm a_f(\bm p)/\|\bm a_f(\bm p)\|_2$。这一操作对同一位置的整列响应除以同一个数，站间相对幅度不变。逐频源系数自由估计时，列缩放可由相应系数补偿，因此不改变联合拟合的最小残差。离散扫描包含站点坐标时，为避免$d_l=0$的数值奇点，实现将不足1\,m的距离按1\,m计算；本文的导数推导适用于这一保护区外，仿真源位置均远离站点。

下文粗体小写和大写分别表示向量与矩阵，$\C$和$\R$分别表示复数域和实数域；$(\cdot)^T$、$(\cdot)^H$及$(\cdot)^\dagger$分别表示转置、共轭转置和Moore--Penrose伪逆，$\|\cdot\|_2$表示欧氏范数，$\bm I_d$表示$d$阶单位矩阵。未标阶数的单位矩阵由相乘矩阵的维度确定。

\subsection{带内信噪比与强弱源条件}
记第$k$个源单独在第$l$站产生的时间平均接收功率为$P_{lk}$，名义信号带宽为$B$。对白噪声，带宽为$B$的频带占采样频带的比例为$B/f_s$，故其带内期望功率为$\sigma_t^2B/f_s$。以各站接收功率的平均值定义带内信噪比，以两个源接收总功率之比定义功率差：
\begin{equation}
 \mathrm{SNR}_{k,\mathrm{in}}
 =10\log_{10}\frac{L^{-1}\sum_{l=1}^L P_{lk}}
 {\sigma_t^2B/f_s},\qquad
 \Delta=10\log_{10}\frac{\sum_l P_{l1}}{\sum_l P_{l2}} .
 \label{eq:snr}
\end{equation}
式中$\mathrm{SNR}_{k,\mathrm{in}}$和$\Delta$均以dB表示。信噪比按名义信号带宽$B$定义，与含保护带的实际处理带宽区分。实验固定源2的波形、功率和噪声，仅缩放源1；这样改变功率差时，弱源自身的SNR保持不变。$\Delta>0$时称源1为强源、源2为弱源；$\Delta=0$时保留这两个身份编号。各站的局部功率差由传播距离决定，不必等于四站平均功率差。

\subsection{未知波形消元与联合位置估计}
定位时将逐频未知波形视为确定性待估系数。根据式\eqref{eq:frequency}的独立、等方差复高斯噪声模型，负对数似然中与位置和波形有关的部分为$\sigma_f^{-2}\sum_f\|\bm x_f-\bm A_f\bm s_f\|_2^2$，其中$\sigma_f^2=\sigma_t^2/N$。因此，最大似然估计等价于联合最小二乘估计\cite{ChenHudsonYao2002}。取观测总能量$E_x=\sum_{f\in\mathcal F}\|\bm x_f\|_2^2>0$作固定归一化，定义
\begin{align}
 \mathcal L(\bm P,\{\bm s_f\})
   &=E_x^{-1}\sum_f\|\bm x_f-\bm A_f(\bm P)\bm s_f\|_2^2,\label{eq:fullLS}\\
 \widehat{\bm s}_f(\bm P)&=\bm A_f(\bm P)^\dagger\bm x_f,\label{eq:waveLS}\\
 J(\bm P)&=E_x^{-1}\sum_f\|\bm\Pi_{\bm A_f}^{\perp}\bm x_f\|_2^2,
 \quad \bm\Pi_{\bm A_f}^{\perp}=\bm I-\bm A_f\bm A_f^\dagger .\label{eq:profile}
\end{align}
这里所有求和均遍历$\mathcal F$。给定一组位置，式\eqref{eq:waveLS}直接求出最优波形；$\bm A_f\bm A_f^\dagger$提取这组位置能够解释的观测，$\bm\Pi_{\bm A_f}^{\perp}\in\C^{L\times L}$保留剩余部分。将最优波形代回式\eqref{eq:fullLS}，就得到只含位置变量的目标$J$，即波形消元\cite{GolubPereyra1973}。$J$为无量纲残差，值越小，位置组合对观测的拟合越好。

所求位置为$\widehat{\bm P}=\arg\min_{\bm P\in\Omega^K}J(\bm P)$。当$K=2$时，非线性搜索只需处理两个二维位置，共四个实数。每次试探新位置都共同重估两源波形，使位置比较始终基于两源的完整拟合。

单源最小二乘谱则为
\begin{equation}
 Q_0(\bm p)=E_x^{-1}\sum_f
 \frac{|\bm a_f(\bm p)^H\bm x_f|^2}{\|\bm a_f(\bm p)\|_2^2} .\label{eq:single}
\end{equation}
式中$|\cdot|$表示复数模。单源情况下有$J(\bm p)=1-Q_0(\bm p)$，故最大化$Q_0$与最小化$J$等价。双源情况下，从同一张单源谱取两个峰没有同时拟合两源，不能代替对位置组合$\bm P$的优化。后续算法以单源谱发现初值，再用多源拟合处理相互影响。
